%This is my super simple Real Analysis Homework template

\documentclass{article}
\usepackage[utf8]{inputenc}
\usepackage[english]{babel}
\usepackage[]{amsthm} %lets us use \begin{proof}
\usepackage{amsmath}
\usepackage[]{amssymb} %gives us the character \varnothing

\title{Homework 4}
\author{Aine Zhang}
\date\today
%This information doesn't actually show up on your document unless you use the maketitle command below

\begin{document}
\maketitle %This command prints the title based on information entered above

%Section and subsection automatically number unless you put the asterisk next to them.
\section*{Exercise 4.6.9., Problem 1}
Let $f : \mathbb{R}^n \to \mathbb{R}$ be a $C^1$ function, and let $S \subset \mathbb{R}^{n+1} $ be the graph of $f$. Fix a point $c \in \mathbb{R}^n$, and let $T$ be the tangent hyperplane to $S$ at the point $C = (c, f (c)) \in S$.
\subsection*{i.}
For any $v \in \mathbb{R}^n$, define $\gamma_v : \mathbb{R} \to \mathbb{R}^2$ by $\gamma_v(t) = tv + c$. Let $\phi_v : \mathbb{R} \to \mathbb{R}^{n+1}$ be defined by $\phi_v(t) = (\gamma_v(t), f (\gamma_v(t)))$. Show that $C + \phi_v'(0)$ lies in $T$.
\begin{proof}
Since $\gamma_v(t) = c + tv,$ we have $\gamma_v'(t) = v$. By the chain rule,
\[
\phi_v'(t)
=
\left(\gamma_v'(t), \frac{d}{dt} f(\gamma_v(t))\right)
=
\left(v, Df(\gamma_v(t))\,v\right).
\]
Evaluating at $t=0$ gives $\phi_v'(0) =(v, Df(c)\,v).$
As for the tangent hyperplane $T$. Since $S$ is the graph of $f$, it can be written as the zero set of $g(x_1,\dots,x_n,x_{n+1}) = f(x_1,\dots,x_n) - x_{n+1}.$
Thus, $S = \{x \in \mathbb{R}^{n+1} : g(x)=0\}.$
The gradient of $g$ is
$$\nabla g(x)
=
\left(D_1 f(x_1,\dots,x_n), \dots, D_n f(x_1,\dots,x_n), -1\right).$$
At the point $C = (c, f(c))$, $\nabla g(C) = (\nabla f(c), -1).$
The tangent hyperplane $T$ at $C$ consists of all vectors $w$ such that $\nabla g(C) \cdot w = 0.$ \\
Now compute $\nabla g(C) \cdot \phi_v'(0) = (\nabla f(c), -1) \cdot (v, Df(c)v).$
This equals $\nabla f(c)\cdot v - Df(c)v.$
Since $Df(c)v = \nabla f(c)\cdot v$, we obtain $\nabla g(C) \cdot \phi_v'(0) = 0.$

Thus $\phi_v'(0)$ is orthogonal to $\nabla g(C)$, so $\phi_v'(0) \in T$. Therefore, $C + \phi_v'(0)$ lies in the tangent hyperplane $T$.
\end{proof}

\subsection*{ii.}
Show that every vector $X \in T$ can be written in the form $X = C + V$, where $V = \phi'(0)$ for some $v \in \mathbb{R}^n$. Hence, every vector in the tangent hyperplane may be realized as the tangent vector to a curve in $S$.

\begin{proof}
Let $X \in T$. Since $T$ is the tangent hyperplane at $C$, we may write $X = C + W$ for some vector $W$ orthogonal to $\nabla g(C)$. Thus $\nabla g(C) \cdot W = 0.$ \\ 
Write $W = (a_1,\dots,a_n,a_{n+1}).$ Set $v := (a_1,\dots,a_n) \in \mathbb{R}^n.$
From part (i) we know that $\phi_v'(0) = (v, Df(c)v).$ \\
Since $W \in T$, we have$\nabla g(C) \cdot W = 0.$
Using $\nabla g(C) = (\nabla f(c), -1)$, this gives
\[
(\nabla f(c), -1)\cdot (v, a_{n+1}) = 0 \Rightarrow
\nabla f(c)\cdot v - a_{n+1} = 0 \Rightarrow a_{n+1} = \nabla f(c)\cdot v.
\]

But $Df(c)v = \nabla f(c)\cdot v$, so $W = (v, Df(c)v) = \phi_v'(0).$ Therefore$ X = C + W = C + \phi_v'(0),$ showing that every vector in $T$ can be written in this form.
Hence, every vector in the tangent hyperplane is realized as the tangent vector at $t=0$ to the curve $\phi_v(t) = (c + tv, f(c+tv))$ in $S$.
\end{proof}

\clearpage %Gives us a page break before the next section. Optional.

\section*{4.5.11., Problem 2}
Let $p_j = (x_j , y_j ), j = 1, ..., m$, be $m$ points in $\mathbb{R}^2$ with at least two $x_j$'s. Given a line $y = mx + b$, define $E(m, b) = \Sigma_{j=1}^n(y_j - (mx_j + b))^2$. Find the values of $m$ and $b$ that minimize this sum. For
the values of $m$ and $b$ that minimize the function $E$, the line $y = mx + b $is called the \textit{ordinary least squares approximation} to the data $p_1, . . . , p_m$.
\begin{proof}
Let $E(m,b)=\sum_{j=1}^n (y_j-(mx_j+b))^2$. Then 
\[
\nabla E(m,b)=\begin{bmatrix}
-2\sum_{j=1}^n (y_j-(mx_j+b))x_j \\
-2\sum_{j=1}^n (y_j-(mx_j+b))
\end{bmatrix}.
\]
Setting $\nabla E=0$ gives the system 
$\sum_{j=1}^n (y_j-(mx_j+b))x_j=0$ and $\sum_{j=1}^n (y_j-(mx_j+b))=0$. 

From the first equation we obtain $\sum_{j=1}^n y_j x_j = m\sum_{j=1}^n x_j^2 + b\sum_{j=1}^n x_j$, and from the second equation we obtain $\sum_{j=1}^n y_j = m\sum_{j=1}^n x_j + nb$. This is a linear system in $m$ and $b$. Solving it yields 
\[
m=\frac{n\sum_{j=1}^n x_j y_j - \left(\sum_{j=1}^n x_j\right)\left(\sum_{j=1}^n y_j\right)}{n\sum_{j=1}^n x_j^2 - \left(\sum_{j=1}^n x_j\right)^2}
\]
and 
\[
b=\frac{\left(\sum_{j=1}^n x_j^2\right)\left(\sum_{j=1}^n y_j\right) - \left(\sum_{j=1}^n x_j\right)\left(\sum_{j=1}^n x_j y_j\right)}{n\sum_{j=1}^n x_j^2 - \left(\sum_{j=1}^n x_j\right)^2}.
\]

Since at least two of the $x_j$ are distinct, the denominator $n\sum_{j=1}^n x_j^2 - (\sum_{j=1}^n x_j)^2$ is nonzero, so the solution is well-defined. Moreover, $E(m,b)\to\infty$ as $\|(m,b)\|\to\infty$, because $E$ is a quadratic polynomial in $m$ and $b$ with positive definite quadratic part. Hence it suffices to examine critical points, and the above solution gives the unique global minimum. Therefore, these values of $m$ and $b$ minimize $E(m,b)$.
\end{proof}

\clearpage

\section*{Exercise 4.7.17., Problem 3}
Let $f : U \subset \mathbb{R}^2 \to \mathbb{R}$ be $C^2$. Let $p \in U$ be a critical point of $f$ and suppose that the matrix of $H_p$ has negative determinant. Show that $f$ does not have a local extremum at $p$. \\
Hint: a 2 × 2 matrix has negative determinant if and only if two eigenvalues are of opposite signs. Study the value of $f$ in the direction of eigenvectors. 
\begin{proof}
Let $p \in U$ be a critical point of $f$, so $\nabla f(p)=0$. Since $f$ is $C^2$, Taylor's theorem gives the expansion
\[
f(p+h)=f(p)+\tfrac12 \langle H_p h, h\rangle + o(\|h\|^2)
\]
as $h \to 0$, where $H_p$ is the Hessian matrix of $f$ at $p$.

We are given that $\det(H_p)<0$. For a $2\times 2$ symmetric matrix, this is equivalent to saying that its two eigenvalues have opposite signs by the hint. Let $\lambda_1>0$ and $\lambda_2<0$ be the eigenvalues of $H_p$, with corresponding eigenvectors $v_1$ and $v_2$.

Consider $h=t v_1$. Since $v_1$ is an eigenvector, we have $H_p v_1=\lambda_1 v_1$, and hence $\langle H_p v_1, v_1\rangle=\lambda_1 \|v_1\|^2>0$. Thus
\[
f(p+t v_1)=f(p)+\tfrac12 t^2 \lambda_1 \|v_1\|^2 + o(t^2).
\]
For sufficiently small $t\neq 0$, the quadratic term dominates, and therefore $f(p+t v_1)>f(p)$.

Now if we consider $h=t v_2$. Similarly, $H_p v_2=\lambda_2 v_2$, so $\langle H_p v_2, v_2\rangle=\lambda_2 \|v_2\|^2<0$. Hence
\[
f(p+t v_2)=f(p)+\tfrac12 t^2 \lambda_2 \|v_2\|^2 + o(t^2),
\]
and for sufficiently small $t\neq 0$, we have $f(p+t v_2)<f(p)$.

Thus, in one direction near $p$ the function increases, and in another direction it decreases. Therefore $p$ cannot be a local maximum or a local minimum. Hence $f$ does not have a local extremum at $p$.
\end{proof}
\clearpage

\section*{Exercise 4.7.18, Problem 4}
 Find the critical points of $f(x, y) = x^3 + 8y^3 - 6xy - 2$. For each, determine if it is a local maximum, local minimum, or saddle point, if possible. 

 \begin{proof}
     We want to find where $\nabla f=0$, so we take the partial derivatives and get $\nabla f =\begin{bmatrix}
         3x^2-6y & 24y^2-6x
     \end{bmatrix}=0. $ So $ 3x^2-6y=0, 24y^2-6x=0 \Rightarrow x^2=2y, x=4y^2 \Rightarrow x^4=4y^2 \Rightarrow x^4=x \Rightarrow x(x^3-1)=0 \Rightarrow x = 0 \text{ or } x=1$. \\
     So, $\nabla f =0 \text{ when }\begin{cases}
         x =0, y=0 \\
         x = 1, y = 0.5
     \end{cases}$ \\
     We want to apply the second derivative test to see what happens to the two points. \\
     $$Hf(x)=\begin{bmatrix}
         D_1^2f & D_1D_2f \\
         D_2D_1f & D_2^2f
     \end{bmatrix}.$$
    We know $D_1^2f=6x, D_2D_1f=D_1D_2f=-6, D_2^2=48y$, so we get\\
         $$Hf(x)=\begin{bmatrix}
         6x  & -6 \\
         -6 & 48y
     \end{bmatrix}.$$
    Plug in $x=1, y=0.5$ now, we get 
        $Hf(x)=\begin{bmatrix}
         6 & -6 \\
         -6 & 24
     \end{bmatrix}.$ Next we find the eigenvalues:         $$Char(\begin{bmatrix}
         6 & -6 \\
         -6 & 24
     \end{bmatrix})= \det(\begin{bmatrix}
         \lambda & 0 \\
         0 & \lambda
     \end{bmatrix} - \begin{bmatrix}
         6 & -6 \\
         -6 & 24 
    \end{bmatrix})= \det(\begin{bmatrix}
         \lambda - 6 & 6 \\
         6 & \lambda -24 
    \end{bmatrix})=\lambda^2-30\lambda +108
         $$
    So we get $\lambda = 15 \pm 3\sqrt{13}, \lambda_1=15 + 3\sqrt{13} \approx 25.81, \lambda_2= 15 - 3\sqrt{13} \approx 4.18$. So we get that both eigenvalues are positive, so the Hessian matrix at this point is positive definite and the point $(1,0.5)$ is a local minimum.  \\
    Next, plug in $x=0, y=0$ first, we get 
          $H(f)(x)=\begin{bmatrix}
         0 & -6 \\
         -6 & 0
     \end{bmatrix}.$ Next we find the eigenvalues           $$Char(\begin{bmatrix}
         0 & -6 \\
         -6 & 0
     \end{bmatrix})= \det(\begin{bmatrix}
         \lambda & 0 \\
         0 & \lambda
     \end{bmatrix} - \begin{bmatrix}
         0 & -6 \\
         -6 & 0
     \end{bmatrix}) = \det(\begin{bmatrix}
         \lambda & 6 \\
         6 & \lambda
     \end{bmatrix}) =\lambda^2-36 \Rightarrow \lambda = \pm 6 $$
     The Hessian matrix is neither positive definite nor negative definite, so we need extra methods to determine what this point is. \\
     $f(0,0)=-2$. Consider the function along the line $y=0$. Then $f(x,0)=x^3-2$. For $x>0$ small, $f(x,0)>-2$, and for $x<0$ small, $f(x,0)<-2$. Thus, in any neighborhood of $(0,0)$, the function takes values both greater than and less than $f(0,0)$. Therefore, $(0,0)$ is a saddle point.

     Hence, the critical points are $(0,0)$, which is a saddle point, and $(1,0.5)$, which is a local minimum.
\end{proof}


\end{document}